% preamble.tex — MathJax-safe shorthands, converted from myshorthand.sty
% (place this at scaffolds/preamble.tex in your Hexo project)

% ---- Paired delimiters ----
% mathtools' \DeclarePairedDelimiter(X) has no MathJax equivalent, so these
% are rewritten as plain \newcommand using \left...\right / \middle, which
% MathJax's core + ams packages support natively.
\newcommand{\abs}[1]{\left\lvert #1 \right\rvert}
\newcommand{\eval}[1]{\left. #1 \right\rvert}
\newcommand{\ket}[1]{\left\lvert #1 \right\rangle}
\newcommand{\bra}[1]{\left\langle #1 \right\rvert}
\newcommand{\braket}[2]{\left\langle #1 \,\middle|\, #2 \right\rangle}
\newcommand{\braHket}[3]{\left\langle #1 \,\middle|\, #2 \,\middle|\, #3 \right\rangle}
\newcommand{\Set}[2]{\left\{ #1 \,\middle|\, #2 \right\}}

% ---- Number sets ----
\newcommand{\reals}{\mathbb{R}}
\newcommand{\rationals}{\mathbb{Q}}
\newcommand{\integers}{\mathbb{Z}}
\newcommand{\naturals}{\mathbb{N}}
\newcommand{\complex}{\mathbb{C}}

% ---- Transforms ----
\newcommand{\fourier}[1]{\mathcal{F}\left\{#1\right\}}
\newcommand{\invfourier}[1]{\mathcal{F}^{-1}\left\{#1\right\}}
\newcommand{\laplace}[1]{\mathcal{L}\left\{#1\right\}}
\newcommand{\invlaplace}[1]{\mathcal{L}^{-1}\left\{#1\right\}}

% ---- Probability / statistics ----
% \ensuremath dropped: MathJax macros always expand inside math mode already,
% so it's a redundant no-op here.
\newcommand{\expectation}[2][]{\mathbb{E}_{#1}\left[#2\right]}
\newcommand{\variance}[2][]{\mathbb{V}_{#1}\left[#2\right]}
\DeclareMathOperator{\covar}{Cov}

% ---- Common functions ----
\DeclareMathOperator{\rect}{rect}
\DeclareMathOperator{\tri}{tri}
\DeclareMathOperator{\sinc}{sinc}
\DeclareMathOperator{\sgn}{sgn}
% \char"0428 (raw font glyph access) isn't supported by MathJax; \unicode{}
% is MathJax's built-in equivalent for inserting an arbitrary code point.
\newcommand{\Sha}{\unicode{x0428}}
\newcommand{\comb}[1]{\Sha_{#1}}

% ---- Operators ----
\DeclareMathOperator{\grad}{grad}
\DeclareMathOperator{\divv}{div}
\DeclareMathOperator{\curl}{curl}
\DeclareMathOperator{\idd}{id}
\DeclareMathOperator{\supp}{supp}
\DeclareMathOperator{\rank}{rank}
\DeclareMathOperator{\adj}{adj}
\DeclareMathOperator{\vol}{vol}
\DeclareMathOperator{\ev}{ev}
\DeclareMathOperator{\tr}{tr}
\newcommand{\iso}{\simeq}
\newcommand{\hodge}{\mathord{\star}}
\newcommand{\codiff}{\mathord{\delta}}
\newcommand{\LdR}{\mathord{\Delta}}
\newcommand{\dAlembert}{\mathord{\Box}}
\newcommand{\iprod}{\mathbin{\lrcorner}}
\newcommand{\indUp}[1]{{#1}^{\sharp}}
\newcommand{\indDown}[1]{{#1}^{\flat}}
\newcommand{\transpose}[1]{{#1}^{\top}}
\newcommand{\hermitian}[1]{{#1}^{\dagger}}

% ---- Miscellaneous symbols and shorthands ----
\renewcommand{\Re}{\mathfrak{Re}}
\renewcommand{\Im}{\mathfrak{Im}}
\newcommand{\powerfactor}{\mathrm{pf}}
\renewcommand{\vec}[1]{\boldsymbol{\mathbf{#1}}}
\newcommand{\lalg}[1]{\mathfrak{#1}}
